% !TEX root = ../jlt.tex
\section{Jlt, by Three Examples}
\label{sec:examples}

This section makes the case in three listings, each on a function the reader already knows.
The first shows a bug that a signature would have prevented.
The second shows an optimisation that a declared law permits.
The third shows a fact that the compiler works out from facts it was given.
Each listing has its code on the left, the tests of that code on the right, and below them what the code becomes under the translation of Section~\ref{sec:translation}.

\subsection{What a signature prevents}
\label{sec:examples:signature}

Listing~\ref{lst:motivation} is about a function every \Cpp{} programmer has written.

\begin{listing}[p]
\begin{minipage}[t]{0.48\linewidth}
\begin{cpp}
template <typename T, typename Op>
constexpr bool associative = false;
template <typename T, typename Op,
          typename E>
constexpr bool unital = false;
template <typename T, typename Op>
concept IsMagma = copyable<T> &&
  requires(T a, T b) {
    { Op{}(a, b) } -> same_as<T>; };
template <typename T, typename Op,
          typename E>
concept IsUnitalMagma =
  IsMagma<T, Op> && unital<T, Op, E> &&
  requires { { E{}() } -> same_as<T>; };
template <typename T, typename Op,
          typename E>
concept IsMonoid =
  IsUnitalMagma<T, Op, E> &&
  associative<T, Op>;
\end{cpp}
\end{minipage}\hfill
\begin{minipage}[t]{0.48\linewidth}
\begin{cpp}
struct True { constexpr bool
  operator()() const { return true; } };
template <> constexpr bool
associative<bool, logical_and<bool>> = true;
template <> constexpr bool
unital<bool, logical_and<bool>, True> = true;
static_assert(logical_and<bool>{}(
    True{}(), True{}()) == True{}());

struct Bare { constexpr bool
  operator==(const Bare&) const = default; };
struct BareOp { constexpr Bare
  operator()(Bare, Bare) const { return {}; } };
struct BareUnit { constexpr Bare
  operator()() const { return {}; } };
template <> constexpr bool
associative<Bare, BareOp> = true;
template <> constexpr bool
unital<Bare, BareOp, BareUnit> = true;
\end{cpp}
\end{minipage}

\begin{minipage}[t]{0.48\linewidth}
\begin{cpp}
template <ranges::range R>
constexpr ranges::range_value_t<R>
concat1(R r) {
  ranges::range_value_t<R> s{};
  for (auto x : r) s = s + x;
  return s; }
\end{cpp}
\end{minipage}\hfill
\begin{minipage}[t]{0.48\linewidth}
\begin{cpp}
template <typename R>
concept Accepts1 =
    requires(R r) { concat1(r); };
static_assert(
concat1(array{1uz, 2uz, 3uz}) == 6 &&
concat1(array{"ab"s, "c"s}) == "abc"s &&
concat1(array{true, false}) == true &&
Accepts1<array<Bare, 2>>);
\end{cpp}
\end{minipage}

\begin{minipage}[t]{0.48\linewidth}
\begin{cpp}
template <typename T, typename Op,
          typename E, ranges::sized_range R>
  requires IsMonoid<T, Op, E> && same_as<
      ranges::range_value_t<R>, T>
constexpr T concat2(R r) {
  return ranges::fold_left(
      r, E{}(), Op{}); }
\end{cpp}
\end{minipage}\hfill
\begin{minipage}[t]{0.48\linewidth}
\begin{cpp}
static_assert(
concat2<size_t, plus<size_t>, Zero>(
    array{1uz, 2uz, 3uz}) == 6 &&
concat2<string, plus<string>, Empty>(
    array{"ab"s, "c"s}) == "abc"s &&
concat2<bool, logical_and<bool>, True>(
    array{true, false}) == false &&
concat2<Bare, BareOp, BareUnit>(
    array{Bare{}, Bare{}}) == Bare{});
\end{cpp}
\end{minipage}
{\small\[
\begin{array}{@{}l@{}}
\mathit{Magma}\;T = \{\,\mathit{op} : T \to T \to T\,\} \\[2pt]
\mathit{UnitalMagma}\;T = \{\,\mathit{magma} : \mathit{Magma}\;T,\ \mathit{unital} : \mathit{Bool},\ \mathit{unit} : T\,\} \\[2pt]
\mathit{Monoid}\;T = \{\,\mathit{um} : \mathit{UnitalMagma}\;T,\ \mathit{assoc} : \mathit{Bool}\,\} \\[4pt]
\mathit{concat2} = \Lambda T.\,\Lambda \mathit{Op}.\,\Lambda E.\,\lambda q : \mathit{Monoid}\;T.\,\lambda n.\,\lambda f. \\
\qquad\qquad\quad \pi_2\,(\mathit{iter}\;n\;(\lambda (k, s).\,(k + 1,\ q.\mathit{um}.\mathit{magma}.\mathit{op}\;s\;(f\;k)))\;(0,\ q.\mathit{um}.\mathit{unit}))
\end{array}
\]}
\caption{What a signature prevents: a function that keeps less than it promises, and one held to its signature.
Left: a monoid as a shape; \cppinline{concat1}; \cppinline{concat2}.
Right: two monoids declared, the second the smallest there is; a test of each function.
Below: the image of \cppinline{concat2} in System~F.}
\label{lst:motivation}
\end{listing}
% TODO(listings): Listing lst:rule2 still shows the earlier form of IsMonoid,
% with the unit as a value, and follows.  The image is the one of the smaller
% kernel (a range as a count and an arrow, the iterator on a pair), which
% Section 4 does not yet describe; it is a translation by hand.  The test that
% expects the compiler to refuse a call of concat1 at Bare is to be written.

The left column is code as it would stand in a program, and the right column tests it.
A monoid is built up as algebra builds it.
A \emph{magma} is a carrier with an operation, a \emph{unital magma} has a unit as well, and a \emph{monoid} is a unital magma whose operation is associative.
The carrier, the operation and the unit are each a type, so the shape reads as the triple it is.
That the operation is associative, and that the unit is its unit, are axioms: the programmer states them, as the top right panel does for the truth values under conjunction, and the compiler takes them at their word.
The monoids of \cppinline{size_t} and of strings are declared the same way, with units \cppinline{Zero} and \cppinline{Empty}.
The unit is also checked at one point, by an assertion beside its axiom, which can refute a false claim and cannot prove a true one.
The shape itself asks of the carrier only that its values can be copied, and in particular no equality, for a reason that Section~\ref{sec:translation} gives.

\cppinline{concat1} is the loop one writes without a second thought: a sum that starts from a default value.
Its signature asks only for a range.
Its body needs more: an operation \cppinline{+} on the elements, and a default value that is that operation's unit.
At \cppinline{size_t} and at strings that is the monoid a caller has in mind.
At the truth values it is disjunction from \cppinline{false}, since the compiler adds two truth values as integers, so \cppinline{concat1} returns \cppinline{true} for a range of \cppinline{true} and \cppinline{false}.
That is right for a caller who meant ``some'' and wrong for one who meant ``all'', and nothing in the signature says which.
\cppinline{Bare} is the smallest monoid there is: a type with one operation and one unit, and an equality only so that the test can state a result.
Asked whether \cppinline{concat1} accepts a range of it, \Cpp{} answers yes, as the last line of the test shows, because it reads the signature and not the body.
The call itself does not compile.
That is an experiment a reader can make on a template of their own: hand it the barest type that meets its constraints.

\cppinline{concat2} names in its signature what it uses, and its body uses nothing else.
So it is right at all three monoids, and it is defined at \cppinline{Bare}.
Its image below the panels says the same in another language: a shape becomes a record with one field for each thing the shape asks, a nested shape a nested record, and the function takes that record as an argument.
A range is there its length $n$ and the arrow $f$ from an index to an element, and the body reaches a value of $T$ only through $q$.
\cppinline{concat1} has no image: no record gives it a \cppinline{+}, and a checker that holds a body to its signature refuses it and says so, that no operation \cppinline{+} is declared for the range's elements.
\Cpp{} accepts both functions.

Of the lines of Listing~\ref{lst:motivation}, the shapes, the traits with their entries and the function objects are in the kernel as written.
The range that \cppinline{concat2} folds over, and the arrays and strings of the tests, are derived forms of it (Section~\ref{sec:kernel:derived}).

\subsection{What a law permits}
\label{sec:examples:law}

The function of Listing~\ref{lst:motivation} is not new.
A Haskell programmer knows it as \texttt{mconcat}, and knows its companion \texttt{stimes}, which combines $n$ copies of one value~\cite{ghc2025semigroup}.
Listing~\ref{lst:power} is that companion, the \emph{power} of a monoid.

\begin{listing}[p]
\begin{minipage}[t]{0.48\linewidth}
\begin{cpp}
template <typename T> struct Power {
  T product, base; size_t exponent; };
template <typename T, typename Op,
          typename E>
struct Halving {
  constexpr Power<T> operator()(
      Power<T> p, size_t) const
    requires IsMonoid<T, Op, E> {
    return {
      p.exponent % 2 == 1
          ? Op{}(p.product, p.base) : p.product,
      Op{}(p.base, p.base),
      p.exponent / 2}; } };
template <typename T, typename Op,
          typename E>
  requires IsMonoid<T, Op, E>
constexpr T power(T x, size_t n) {
  return ranges::fold_left(
      views::iota(0uz, size_t(bit_width(n))),
      Power<T>{E{}(), x, n},
      Halving<T, Op, E>{}).product; }
\end{cpp}
\end{minipage}\hfill
\begin{minipage}[t]{0.48\linewidth}
\begin{cpp}
static_assert(
power<size_t, plus<size_t>, Zero>(
    2uz, 10uz) ==
concat2<size_t, plus<size_t>, Zero>(
    views::repeat(2uz, 10uz)) &&
power<size_t, plus<size_t>, Zero>(
    2uz, 10uz) == 20 &&
power<size_t, multiplies<size_t>, One>(
    2uz, 10uz) == 1024 &&
power<string, plus<string>, Empty>(
    "ab"s, 3uz) == "ababab"s &&
power<size_t, multiplies<size_t>, One>(
    3uz, 1'000'000'000uz) ==
    9'215'150'800'219'179'009uz &&
power<string, plus<string>, Empty>(
    ""s, 1'000'000'000'000uz) == ""s);
\end{cpp}
\end{minipage}

\begin{minipage}[t]{0.48\linewidth}
\begin{cpp}
template <typename T, typename Op,
          typename E>
  requires IsIdempotentMonoid<T, Op, E>
constexpr T power(T x, size_t n) {
  return n == 0 ? E{}() : x; }
template <typename T, typename Op,
          typename E>
  requires IsMonoid<T, Op, E>
constexpr T guessed_power(T x, size_t n) {
  return power<T, Op, E>(x, n); }
\end{cpp}
\end{minipage}\hfill
\begin{minipage}[t]{0.48\linewidth}
\begin{cpp}
static_assert(
guessed_power<size_t, plus<size_t>,
    Zero>(2uz, 10uz) == 20 &&
guessed_power<size_t, Plus2,
    Zero>(2uz, 10uz) == 2);
\end{cpp}
\end{minipage}
{\small\[
\begin{array}{@{}l@{}}
\mathit{power}_0 = \Lambda T.\,\Lambda \mathit{Op}.\,\Lambda E.\,\lambda q : \mathit{Monoid}\;T.\,\lambda x.\,\lambda n.\ \mathit{concat2}\;T\;\mathit{Op}\;E\;q\;n\;(\lambda k.\,x) \\[3pt]
\mathit{power} = \Lambda T.\,\Lambda \mathit{Op}.\,\Lambda E.\,\lambda q : \mathit{Monoid}\;T.\,\lambda x.\,\lambda n.\ \pi_1\,(\mathit{iter}\;(\mathit{bits}\;n)\;\mathit{step}\;(q.\mathit{um}.\mathit{unit},\ x,\ n)) \\
\qquad \mathit{step}\,(p, b, e) = (\,\mathit{if}\ \mathit{odd}\;e\ \mathit{then}\ p \cdot b\ \mathit{else}\ p,\ \ b \cdot b,\ \ e / 2\,) \qquad a \cdot b = q.\mathit{um}.\mathit{magma}.\mathit{op}\;a\;b
\end{array}
\]}
\caption{What a law permits: the power of a monoid, and a caller that names too little.
Left: the power by repeated squaring; a case of it that a law selects, and a function that calls it.
Right: their tests.
Below: the power's definition and the power in System~F.}
\label{lst:power}
\end{listing}
% TODO(listings): code from the probe (namespace triple_probe).  The shapes
% IsCommutativeMonoid (a monoid whose operation is declared commutative) and
% the law commutative are omitted from the panels.  The free theorem as a
% guarded function (fold_commutes, in the probe) is a candidate for an
% appendix.  The citations for the
% lemma and for the Russian peasant method are from memory, to be verified.

The power of $x$ is $x$ combined with itself $n$ times: the concatenation of a repetition, which the first line of the test says in \Cpp{}.
Computed that way it takes $n$ steps.
\cppinline{power} takes as many steps as $n$ has bits, by squaring: the method known for multiplication as that of the Russian peasant~\cite[\S 4.6.3]{knuth1997art2}.
The two agree because the operation is associative, which the shape in the signature states.
It is one function with a meaning at each monoid: a multiple under \cppinline{+}, a power under \cppinline{*}, a repetition at strings.
The images below the panels differ in one place, the count of the iterator, $n$ in the definition and the bits of $n$ in \cppinline{power}.
System~F does not know that they are equal.
The field of the record that says the operation is associative is what licenses putting the second for the first.
The power of the unit is the unit, whatever the count, as the last line of the test shows at strings.
\cppinline{power} has no case for that.
With the base fixed to the unit and the count left for later, the compiler removes the fold and the count by itself, since the unit squared is the unit and it finds that by computing it (Section~\ref{sec:binding}).

The lines of Listing~\ref{lst:power} are in the kernel as written, the fold over a count being its one recursion over a value, apart from the range in the first line of the test (Section~\ref{sec:kernel:derived}).

The lower panels show what the rules are for.
Where the operation is idempotent, as conjunction is, a power is the value itself, and a second case of \cppinline{power}, on the same signature with one more law, says so in one step.
The compiler selects it wherever that law is declared.
\cppinline{guessed_power} declares a monoid and calls \cppinline{power}, so which case it gets is decided where it is used, by a law its signature never named.
\cppinline{Plus2} is addition under another name, which someone has declared idempotent.
That axiom is false, the caller never chose to depend on it, and ten twos come out as two.
\Cpp{} accepts this function, and the rules of this paper do not: a generic function names every law its result depends on (\RuleTwo{}).

\subsection{What the compiler infers}
\label{sec:examples:inference}

Both listings rest on laws that someone declared.
Listing~\ref{lst:theorems} shows the compiler deriving one, and what a derived law is good for.

\begin{listing}[p]
\begin{minipage}[t]{0.48\linewidth}
\begin{cpp}
template <typename T, typename Op,
          typename E>
struct Mon {};
using Strings =
    Mon<string, plus<string>, Empty>;
using Counts =
    Mon<size_t, plus<size_t>, Zero>;
using All =
    Mon<bool, logical_and<bool>, True>;
using Any =
    Mon<bool, logical_or<bool>, False>;
template <typename H, typename A,
          typename B>
constexpr bool homomorphism = false;
\end{cpp}
\end{minipage}\hfill
\begin{minipage}[t]{0.48\linewidth}
\begin{cpp}
struct Size { constexpr size_t
  operator()(string s) const {
    return s.size(); } };
struct IsEmpty { constexpr bool
  operator()(string s) const {
    return s.empty(); } };
struct Not { constexpr bool
  operator()(bool b) const { return !b; } };
template <> constexpr bool
homomorphism<Size, Strings, Counts> = true;
template <> constexpr bool
homomorphism<IsEmpty, Strings, All> = true;
template <> constexpr bool
homomorphism<Not, All, Any> = true;
\end{cpp}
\end{minipage}

\vspace{-7pt}
\begin{minipage}[t]{0.48\linewidth}
\begin{cpp}
template <typename F, typename B,
    typename G, typename X,
    typename Y, typename Z>
struct Via { F f; G g;
  constexpr Z operator()(X x) const
    requires Maps<F, X, Y> &&
             Maps<G, Y, Z> {
    return g(f(x)); } };
template <typename F, typename B,
    typename G, typename X,
    typename Y, typename Z,
    typename A, typename C>
constexpr bool homomorphism<
    Via<F, B, G, X, Y, Z>, A, C> =
  homomorphism<F, A, B> &&
  homomorphism<G, B, C>;
\end{cpp}
\end{minipage}\hfill
\begin{minipage}[t]{0.48\linewidth}
\begin{cpp}
using NonEmpty = Via<IsEmpty, All, Not,
                     string, bool, bool>;
static_assert(
homomorphism<NonEmpty, Strings, Any> &&
!homomorphism<NonEmpty, Strings, All>);
template <typename H, typename T,
    typename Op, typename E, typename U,
    typename Op2, typename E2>
concept IsHomomorphism =
  IsMonoid<T, Op, E> &&
  IsMonoid<U, Op2, E2> &&
  homomorphism<H, Mon<T, Op, E>,
               Mon<U, Op2, E2>> &&
  requires(H h, T x) {
    { h(x) } -> same_as<U>; };
\end{cpp}
\end{minipage}

\vspace{-7pt}
\begin{minipage}[t]{0.48\linewidth}
\begin{cpp}
template <typename T, typename Op,
    typename E, typename U, typename Op2,
    typename E2, typename H,
    ranges::sized_range R>
  requires IsHomomorphism<H, T, Op, E,
                          U, Op2, E2> &&
    equality_comparable<U> &&
    same_as<ranges::range_value_t<R>, T>
constexpr bool free_theorem(
    Mon<T, Op, E>, Mon<U, Op2, E2>,
    H h, R r) {
  return h(concat2<T, Op, E>(r)) ==
      concat2<U, Op2, E2>(
          views::transform(r, h)); }
\end{cpp}
\end{minipage}\hfill
\begin{minipage}[t]{0.48\linewidth}
\begin{cpp}
static_assert(
free_theorem(Strings{}, Counts{}, Size{},
    array{"ab"s, ""s, "c"s}) &&
free_theorem(Strings{}, All{}, IsEmpty{},
    array{"ab"s, ""s}) &&
free_theorem(Strings{}, Any{}, NonEmpty{},
    array{""s, "ab"s}));
\end{cpp}
\end{minipage}
\vspace{-7pt}
{\small\[
\begin{array}{@{}l@{}}
\mathit{free\_theorem} = \Lambda T.\,\Lambda U.\,\ldots\,\lambda q : \mathit{Monoid}\;T.\,\lambda q' : \mathit{Monoid}\;U.\,\lambda l : \mathit{Bool}.\,\lambda \mathit{eq} : U \to U \to \mathit{Bool}.\,\lambda h : T \to U.\,\lambda n.\,\lambda f. \\
\qquad\qquad \mathit{eq}\;(h\,(\mathit{concat2}\;q\;n\;f))\;(\mathit{concat2}\;q'\;n\;(\lambda k.\,h\,(f\,k)))
\end{array}
\]}
\caption{What the compiler infers: a law derived from laws, and a theorem that comes for free.
Left: the law, a lemma, the theorem. Right: three axioms, the inferred fact, three instances.}
\label{lst:theorems}
\end{listing}
% TODO(listings): code from the probe (namespace triple_probe), green on CI.
% The shape Maps (a required call from one carrier to another) is omitted from
% the panels.
% Parked: the semiring, under which a power is one product, and the lemma that
% a power is a homomorphism where the operation commutes; their code is in the
% probe and their text in this file's history.

A \emph{homomorphism} is a function from one monoid to another that respects their operations and their units.
The listing names each monoid by one type and makes ``is a homomorphism'' a law, as associativity is.
Three axioms are declared: the length of a concatenation is the sum of the lengths, a concatenation is empty exactly when every part is, and De Morgan's law, that the negation of a conjunction is the disjunction of the negations.
One lemma is stated and not proved, since it is a textbook fact: a homomorphism after a homomorphism is a homomorphism.
It is an entry at a composite that reads the same law at the two parts, and the composite names the monoid in the middle, which the compiler could not guess.
\cppinline{Via} also names the three carriers it passes through, and \cppinline{Maps<F, X, Y>} is the shape that asks for a call of \cppinline{F} from \cppinline{X} to \cppinline{Y}.

Nobody declared anything about a string not being empty.
The compiler concludes, from emptiness and De Morgan's law, that it is a homomorphism into the truth values under disjunction: a concatenation is not empty exactly when some part is not.
It concludes nothing more than follows, and the same function against conjunction is refused.

A derived law is worth having because of a theorem that comes with every function written under the rules of this paper: it commutes with every homomorphism of the shape it names, which follows from its type alone~\cite{reynolds1983types,wadler1989theorems}.
\cppinline{free_theorem} states that for \cppinline{concat2} in plain \Cpp{}, with the theorem's hypothesis as its signature: a shape that names the law and the one call the body makes of the function.
Each of its three calls reads as a sentence: from this monoid, to that one, by this function, on these values.
The third can be written only because of the fact the compiler derived.
A call checks one case and proves nothing, and the theorem holds whether or not anyone calls it.
The lines of Listing~\ref{lst:theorems} are in the kernel as written, apart from the ranges (Section~\ref{sec:kernel:derived}).
An axiom, in turn, is only as good as its author: ``a sum of counts is zero exactly when every count is'' looks as safe as the three above, and is false of \cppinline{size_t}, where a sum can wrap around to zero.

\subsection{What the three have in common}
\label{sec:examples:common}
\label{sec:posture}

None of the three listings needs anything but standard \Cpp{}.
What they need is that their signatures can be trusted, and nothing in \Cpp{} ensures that.
What the functions that behave have in common can be said in three rules, which together are the \emph{Juliet Posture}.
Table~\ref{tab:posture} sets them side by side.

\begin{table}[tb]
\centering
\small
\begin{tabularx}{\linewidth}{@{}l>{\raggedright\arraybackslash}X>{\raggedright\arraybackslash}X>{\raggedright\arraybackslash}X>{\raggedright\arraybackslash}X@{}}
\toprule
rule & keeps out of Jlt & buys & costs & held by \\
\midrule
1 Substitution & state, pointers, exceptions & one program, at compile time and at run time & no state beyond a call's own locals, no I/O & $\ast$~in part\newline $\dagger$~syntax \\
2 Abstraction & a body using what its constraints do not name & parametricity & no ad hoc use of a type & $\dagger$~names\newline $\ddagger$~the rest \\
3 Descent & unbounded loops, general recursion & termination & a bound on every fold, stated before it starts & $\ast$~compile time\newline $\dagger$, $\ddagger$~loops \\
\bottomrule
\end{tabularx}
\caption{The posture at a glance: what each rule keeps out of Jlt, what it buys, what it costs and who holds it: $\ast$~the compiler, $\dagger$~the linter, $\ddagger$~review, by a person or an agent, which checks nothing.}
\label{tab:posture}
\end{table}

\paragraph{Rule 1, Substitution.}
\phantomsection\label{def:rule1}
A value is data that can be copied, and a function depends on its arguments alone and can be evaluated while compiling.
Every function in the three listings is \cppinline{constexpr} and takes and returns values, which is why each can be tested by a \cppinline{static_assert}.
The rule keeps out state that outlives a call, pointers, exceptions and input and output.
What it buys is that a value may be replaced by its equal anywhere, and that a program means the same while compiling and while running.

\paragraph{Rule 2, Abstraction.}
\phantomsection\label{def:rule2}
Generic code varies with a type only through what its declaration names.
\cppinline{concat1} breaks it by using a \cppinline{+} that its signature never asked for, and \cppinline{guessed_power} by depending on a law that its signature never named.
What a declaration may name is a \emph{shape}, written as a \cppinline{concept}, and a shape asks for calls and for laws.
A law is a \emph{trait}: a value that the compiler looks up by type, \cppinline{false} unless someone has said otherwise.
An \emph{axiom} is an entry of a trait at one type, stated by the programmer and not proved.
A \emph{lemma} is an entry at a form of type, which reads the same trait at the parts, as the entry at \cppinline{Via} does in Listing~\ref{lst:theorems}.
A \emph{theorem} is what the compiler concludes from the two.
An entry holds for the whole program or not at all, so it is declared above every use that could read it (Appendix~\ref{app:fjlt:order}).

\paragraph{Rule 3, Descent.}
\phantomsection\label{def:rule3}
Recursion is structural, and its number of steps is bounded before the first step.
It has two forms.
Over a value it is the fold of \cppinline{concat2} and of \cppinline{power}:
\begin{equation}
\label{eq:fold}
\begin{aligned}
\mathit{fold}_h\,(e, \langle\rangle) &= e \\
\mathit{fold}_h\,(e, s \cdot r) &= \mathit{fold}_h\,(h(e, s), r) \qquad |r| < |s \cdot r|
\end{aligned}
\end{equation}
where $s \cdot r$ is a range with first element $s$ and tail $r$, $h$ the step and $e$ the start.
The bound is the length of the range: $n$ for \cppinline{concat2}, and the number of bits of $n$ for \cppinline{power}.
Over a type it is the lemma of Listing~\ref{lst:theorems}:
\begin{equation}
\label{eq:cata}
\begin{aligned}
g(x) &= d_x \\
g(t_1 \gg t_2) &= g(t_1) \otimes g(t_2) \qquad t_1, t_2 \triangleleft t_1 \gg t_2
\end{aligned}
\end{equation}
where $x$ is a type with no parts, $d_x$ what is declared for it, $t_1 \gg t_2$ a composite, and $t' \triangleleft t$ says that $t'$ is a proper part of $t$.
A proper part is smaller than its whole, so the compiler's search ends.
The rule keeps out a loop with no stated end and a function that calls itself.

The three rules are a heuristic: what a programmer keeps in mind.
They do not define a language, and keeping them does not by itself put a program inside one.
The rest of the paper supplies that: a small language in which a signature means what it says (Section~\ref{sec:kernel}), and a check that a program is written in it (Section~\ref{sec:check}).
